\section{Saran}

Berdasarkan penelitian yang telah dilakukan, terdapat beberapa saran yang dapat dipertimbangkan untuk pengembangan penelitian selanjutnya sebagai berikut.

\begin{enumerate}

\item Penelitian ini menggunakan bentuk \textit{ansatz} deret Frobenius yang memenuhi kondisi batas pada seluruh titik singular yang relevan. Penelitian selanjutnya dapat mengkaji formulasi \textit{ansatz} alternatif, misalnya dengan menerapkan batasan tambahan pada interval di antara dua titik singular. Pendekatan tersebut diharapkan dapat memperluas formulasi matematis sekaligus meningkatkan efisiensi dalam memperoleh solusi numerik.

\item Perhitungan spektrum frekuensi pada penelitian ini dilakukan menggunakan metode pecahan berlanjut (\textit{continued fraction}) tanpa menerapkan teknik peningkatan konvergensi (\textit{convergence improvement}). Akibatnya, pada beberapa kombinasi parameter masih terdapat frekuensi yang belum berhasil diperoleh karena keterbatasan konvergensi numerik. Oleh karena itu, penelitian selanjutnya dapat menerapkan teknik peningkatan konvergensi maupun metode numerik lainnya untuk meningkatkan akurasi perhitungan serta memperluas rentang parameter yang dapat dianalisis.

\item Penelitian selanjutnya dapat memperluas kajian terhadap perturbasi dan mode kuasinormal pada ruang-waktu Schwarzschild de-Sitter \parencite{MossNorman2002}. Selain perturbasi medan skalar ($s=0$), analisis juga dapat dikembangkan pada perturbasi elektromagnetik ($s=1$) dan perturbasi gravitasional ($s=2$), sehingga karakteristik spektrum frekuensi dan kestabilan lubang hitam untuk berbagai jenis perturbasi dapat dibandingkan secara lebih komprehensif.

\end{enumerate}